% Part: computability
% Chapter: recursive-functions
% Section: halting-problem

\documentclass[../../../include/open-logic-section]{subfiles}

\begin{document}

\olfileid[hi]{cmp}{rec}{hlt}
\olsection{विराम समस्या}

सामान्यतः \emph{विराम समस्या} (halting problem) यह निर्णय करने की
समस्या है कि किसी संगणनीय फलन का विनिर्देश~$e$ (मसलन कोई प्रोग्राम) और
संख्या~$n$ दिए होने पर निवेश~$n$ पर उस फलन का संगणन रुकता है या नहीं,
अर्थात् कोई परिणाम देता है या नहीं। प्रसिद्ध रूप से, ऐलन ट्यूरिंग ने
सिद्ध किया कि स्वयं इस समस्या को किसी संगणनीय फलन से हल नहीं किया जा
सकता; अर्थात् फलन
\[
h(e, n) =
\begin{cases}
1 & \text{यदि संगणन $e$, निवेश $n$ पर रुकता है}\\
0 & \text{अन्यथा,}
\end{cases}
\]
संगणनीय नहीं है।

आंशिक पुनरावर्ती फलनों के संदर्भ में, किसी प्रोग्राम के विनिर्देश की
भूमिका क्लिनी के प्रसामान्य-रूप प्रमेय में दिया अनुक्रमांक~$e$ निभा सकता
है। यदि $f$ कोई आंशिक पुनरावर्ती फलन है, तो हर ऐसा~$e$ जिसके लिए
प्रसामान्य-रूप प्रमेय का समीकरण मान्य हो,~$f$ का अनुक्रमांक है। संख्या~$e$
दिए होने पर प्रसामान्य-रूप प्रमेय कहता है कि
\[
\cfind{e}(x) \simeq U(\mu s \; T(e, x, s))
\]
आंशिक पुनरावर्ती है; और प्रत्येक आंशिक पुनरावर्ती $f\colon \Nat
\to \Nat$ के लिए कोई $e \in \Nat$ ऐसा है कि $\cfind{e}(x) \simeq
f(x)$ सभी~$x \in \Nat$ के लिए। वास्तव में, प्रत्येक ऐसे $f$ के लिए केवल
एक नहीं, बल्कि अनंततः अनेक ऐसे~$e$ हैं। \emph{विराम फलन}~$h$ को इस
प्रकार परिभाषित किया जाता है:
\[
h(e, x) =
\begin{cases}
1 & \text{यदि $\cfind{e}(x) \fdefined$}\\
0 & \text{अन्यथा।}
\end{cases}
\]
ध्यान दें कि $h(e, x) = 0$ होता है यदि $\cfind{e}(x) \fundefined$; और ऐसा
तब भी है जब~$e$ किसी आंशिक पुनरावर्ती फलन का अनुक्रमांक है ही नहीं।

\begin{thm}
\ollabel{thm:halting-problem}
विराम फलन~$h$ आंशिक पुनरावर्ती नहीं है।
\end{thm}

\begin{proof}
यदि $h$ आंशिक पुनरावर्ती होता, तो हम परिभाषित कर सकते थे:
\[
d(y) =
\begin{cases}
1 & \text{यदि $h(y, y) = 0$}\\
\umin{x}{x \neq x} & \text{अन्यथा।}
\end{cases}
\]
चूँकि कोई संख्या $x$, $x \neq x$ को संतुष्ट नहीं करती, इसलिए
$\umin{x}{x \neq
x}$ अस्तित्व में नहीं है; अतः $d(y) \fundefined$ तभी
और केवल तभी जब $h(y,y) \neq 0$। इस परिभाषा से निकलता है:
\begin{enumerate}
\item $d(y) \fdefined$ तभी और केवल तभी जब $\cfind{y}(y) \fundefined$
  अथवा $y$ किसी आंशिक पुनरावर्ती फलन का अनुक्रमांक नहीं है।
\item $d(y) \fundefined$ तभी और केवल तभी जब $\cfind{y}(y) \fdefined$।
\end{enumerate}
यदि $h$ आंशिक पुनरावर्ती होता, तो $d$ भी आंशिक पुनरावर्ती होता। अतः
क्लिनी के प्रसामान्य-रूप प्रमेय से उसका कोई अनुक्रमांक~$e_d$ है।
$h(e_d, e_d)$ के मान पर विचार करें। दो संभव प्रकरण हैं: $0$ और~$1$।
\begin{enumerate}
\item यदि $h(e_d, e_d) = 1$, तो $\cfind{e_d}(e_d) \fdefined$। किंतु
  $\cfind{e_d} \simeq d$, और $d(e_d)$ तभी और केवल तभी परिभाषित है जब
  $h(e_d, e_d) = 0$। अतः $h(e_d, e_d) \neq 1$।
\item यदि $h(e_d, e_d) = 0$, तो या तो $e_d$ किसी आंशिक पुनरावर्ती फलन
  का अनुक्रमांक नहीं है, या है और $\cfind{e_d}(e_d)
  \fundefined$। किंतु
  फिर, $\cfind{e_d} \simeq d$, और $d(e_d)$ तभी और केवल तभी अपरिभाषित है
  जब $\cfind{e_d}(e_d) \fdefined$।
\end{enumerate}
निष्कर्ष यह है कि अंततः $e_d$ किसी आंशिक पुनरावर्ती फलन का अनुक्रमांक
नहीं हो सकता। किंतु यदि $h$ आंशिक पुनरावर्ती होता, तो $d$ भी ऐसा होता;
अतः उसके अनुक्रमांक के रूप में $e_d$ की हमारी परिभाषा मान्य होती। हमें
निष्कर्ष निकालना होगा कि $h$ आंशिक पुनरावर्ती नहीं हो सकता।
\end{proof}

\end{document}

